% 2.5 主动与被动算符；2.6 表示的对称化与反对称化乘积
\section{主动与被动算符}

为完整起见，这里说明由主动操作构成的群的表示的基函数与相应被动操作群的表示之间的关系。

设 $X, Y, Z, \ldots$ 是一群算符（例如点群操作），按主动形式考虑时有 $X_{a}Y_{a} = Z_{a}$。再设已给定这些主动操作的一个表示 $\Gamma_{a}^{i}$，使得 $\mathbf{D}_{a}^{i}(X_{a})\mathbf{D}_{a}^{i}(Y_{a}) = \mathbf{D}_{a}^{i}(Z_{a})$，基函数记为 $\langle \phi_{t}^{i}|$，例如

\begin{equation*}
\bar{X}_{a}\phi_{t}^{i} = \sum_{j} \phi_{j}^{i}\mathbf{D}_{a}^{i}(X_{a})_{jt},
\tag{2.5.1}
\end{equation*}

其中 $\bar{X}_{a}$ 是主动函数空间算符（见 1.5 节）。现在若想使用被动形式的算符：首先 $X_{p} = X_{a}^{-1}$，故群的乘法表有所不同，例如此时有 $Y_{p}X_{p} = Z_{p}$。群当然是同一个群；乘法不同是因为把元素映到其逆的映射是一个反自同构。

\textbf{定理 2.5.1.} \textit{若 $\langle \phi_{t}^{i}|$ 是群 $\mathbf{G}$ 的元素 $X_{a}, Y_{a}, Z_{a}, \ldots$ 的表示 $\Gamma_{a}^{i}$ 的基，则令 $\mathbf{D}_{p}^{i}(X_{p}) = \mathbf{D}_{a}^{i}(X_{a}^{-1})$，$\langle \phi_{t}^{i}|$ 也是元素 $X_{p} = X_{a}^{-1}, Y_{p} = Y_{a}^{-1}, Z_{p} = Z_{a}^{-1}, \ldots$ 的表示 $\Gamma_{p}^{i}$ 的基。}

这是因为

\begin{equation*}
\bar{X}_{p}\phi_{t}^{i} = \bar{X}_{a}^{-1}\phi_{t}^{i} = \sum_{j} \phi_{j}^{i}\mathbf{D}_{a}^{i}(X_{a}^{-1})_{jt} = \sum_{j} \phi_{j}^{i}\mathbf{D}_{p}^{i}(X_{p})_{jt}
\tag{2.5.2}
\end{equation*}

且

\begin{equation*}
\begin{aligned}
\mathbf{D}_{p}^{i}(Y_{p})\mathbf{D}_{p}^{i}(X_{p}) &= \mathbf{D}_{a}^{i}(Y_{a}^{-1})\mathbf{D}_{a}^{i}(X_{a}^{-1}) = \mathbf{D}_{a}^{i}(Y_{a}^{-1}X_{a}^{-1}) \\
&= \mathbf{D}_{a}^{i}(Z_{a}^{-1}) = \mathbf{D}_{p}^{i}(Z_{p}).
\end{aligned}
\tag{2.5.3}
\end{equation*}

作为这一定理应用的一个例子，考虑本章给出的那些表。本章使用主动算符，并对完整列出的各个矩阵表示给出了基函数。若使用被动算符，则只要把列出的矩阵换成它们的逆，同样的基函数仍然可用。也就是说，我们写 $X$（意为 $X_{a}$）之处你写 $X$（意为 $X_{p}$）；我们的乘法表含 $XY = Z$（意为 $X_{a}Y_{a} = Z_{a}$）之处，你的乘法表含 $YX = Z$（意为 $Y_{p}X_{p} = Z_{p}$）；$\langle \phi_{t}^{i}|$ 是表示 $\Gamma_{p}^{i}$ 的基且 $\mathbf{D}^{i}(Y)\mathbf{D}^{i}(X) = \mathbf{D}^{i}(Z)$（意为 $\mathbf{D}_{p}^{i}(Y_{p})\mathbf{D}_{p}^{i}(X_{p}) = \mathbf{D}_{p}^{i}(Z_{p})$），而我们为表示 $\Gamma_{a}^{i}$ 写下的正是同一组基 $\langle \phi_{t}^{i}|$ 与 $\mathbf{D}^{i}(X)\mathbf{D}^{i}(Y) = \mathbf{D}^{i}(Z)$（意为 $\mathbf{D}_{a}^{i}(X_{a})\mathbf{D}_{a}^{i}(Y_{a}) = \mathbf{D}_{a}^{i}(Z_{a})$）。$\Gamma_{p}^{i}$ 与 $\Gamma_{a}^{i}$ 的矩阵之间的关系是 $\mathbf{D}_{p}^{i}(X) = \mathbf{D}_{a}^{i}(X^{-1}) = \{\mathbf{D}_{a}^{i}(X)\}^{-1}$。最后提醒一句：若 $\mathbf{D}_{a}^{i}(X_{a})$ 与 $\mathbf{D}_{p}^{i}(X_{p})$ 具有不同的特征标，则尽管基函数相同，以 $i$ 标记的表示 $\Gamma_{p}^{i}$ 与 $\Gamma_{a}^{i}$ 在特征标表中将有不同的标记。

\section{点群表示的对称化与反对称化乘积}

作为本章的结束，简要叙述点群表示的对称化与反对称化乘积。这些乘积本身就有重要意义，例如与分子振动（Jahn and Teller 1937）以及二级相变的 Landau 理论（例如见 Birman (1966\textsuperscript{b})，Indenbom (1960\textsuperscript{a})，Landau and Lifshitz (1958\textsuperscript{b})，Lyubarskii (1960)）有关；此外，研究点群表示的这些乘积，也是 4.8 节空间群表示的对称化与反对称化乘积工作的一个简单引论。

量子力学的著名结果是：一组全同粒子的波函数，对交换其中任意一对全同粒子必须是对称的或反对称的（见任何一本可靠的量子力学教科书）。总波函数对玻色子（例如光子、$\alpha$ 粒子、声子、磁振子等）必须对称，对费米子（例如电子、质子、中子等）必须反对称；费米子的反对称原理更常见的表述是泡利不相容原理。当有两个全同粒子时，这就引出表示的对称化与反对称化克罗内克平方的概念；群的两个表示的克罗内克积已在 1.3 节引入（见式 (1.3.23)）。

设 $\mathbf{S}$ 是一个群，$\Delta$ 是 $\mathbf{S}$ 的表示，基为 $\langle \phi| = \langle \phi_{1}, \phi_{2}, \ldots, \phi_{d}|$，使得对一切 $s \in \mathbf{S}$ 有

\begin{equation*}
s\phi_{i} = \sum_{j=1}^{d} \phi_{j}\Delta(s)_{ji}.
\tag{2.6.1}
\end{equation*}

记表示 $\Delta$ 的载体空间为 $V$，则可考虑 $V$ 与自身的张量积 $V \otimes V$，即由有序函数对 $(\phi_{i}, \phi_{j})$（$i = 1$ 到 $d$，$j = 1$ 到 $d$）的线性组合张成的向量空间。$V \otimes V$ 的维数为 $d^{2}$，它在外直积群 $\mathbf{S} \otimes \mathbf{S}$ 下不变。由于内直积 $\mathbf{S} \boxtimes \mathbf{S}$ 是 $\mathbf{S} \otimes \mathbf{S}$ 的子群，且在自然映射 $(s, s) \rightarrow s$ 下与 $\mathbf{S}$ 同构，因此可把 $V \otimes V$ 视为在 $\mathbf{S}$ 自身下不变的空间，所定义的 $\mathbf{S}$ 的表示称为 $\Delta$ 的\textit{内克罗内克平方}（inner Kronecker square）。这与 1.3 节的思想和记号一致，记此表示为 $\Delta \boxtimes \Delta$。为得到 $\Delta \boxtimes \Delta$ 的显式形式，有

\begin{equation*}
s(\phi_{i}, \phi_{j}) = (s\phi_{i}, s\phi_{j})
\tag{2.6.2}
\end{equation*}

\begin{equation*}
= \sum_{k,l} (\phi_{k}, \phi_{l})\Delta(s)_{ki}\Delta(s)_{lj}
\tag{2.6.3}
\end{equation*}

\begin{equation*}
= \sum_{k,l} (\phi_{k}, \phi_{l})(\Delta \boxtimes \Delta)(s)_{kl,ij}.
\tag{2.6.4}
\end{equation*}

因此

\begin{equation*}
(\Delta \boxtimes \Delta)(s) = \Delta(s) \otimes \Delta(s),
\tag{2.6.5}
\end{equation*}

即矩阵 $\Delta(s)$ 的克罗内克平方；并且由式 (2.6.3)，$\Delta \boxtimes \Delta$ 的特征标为

\begin{equation*}
\chi_{\Delta \boxtimes \Delta}(s) = \sum_{k,l} \Delta(s)_{kk}\Delta(s)_{ll} = \chi_{\Delta}^{2}(s).
\tag{2.6.6}
\end{equation*}

一般地，即使 $\Delta$ 本身不可约，表示 $\Delta \boxtimes \Delta$ 在 $\mathbf{S}$ 上也是可约的。当 $d > 1$ 时，立即可作部分约化：把 $V \otimes V$ 分解为对称部分与反对称部分（例如见 Lyubarskii (1960) 第 4 章）。分别记这两个 $V \otimes V$ 的子空间为 $[V \otimes V]$ 与 $\{V \otimes V\}$，定义如下：$[V \otimes V]$ 是维数为 $\tfrac{1}{2}d(d+1)$ 的向量空间，由对称化对 $(\phi_{i}, \phi_{j}) + (\phi_{j}, \phi_{i})$ 的线性组合张成；$\{V \otimes V\}$ 是维数为 $\tfrac{1}{2}d(d-1)$ 的向量空间，由反对称化对 $(\phi_{i}, \phi_{j}) - (\phi_{j}, \phi_{i})$ 的线性组合张成。为看清这些子空间在 $\mathbf{S}$ 下不变，考虑对称化空间 $[V \otimes V]$：

\begin{equation*}
s\left[(\phi_{i}, \phi_{j}) + (\phi_{j}, \phi_{i})\right] = \sum_{k,l} (\phi_{k}, \phi_{l})\left\{\Delta(s)_{ki}\Delta(s)_{lj} + \Delta(s)_{kj}\Delta(s)_{li}\right\}
\tag{2.6.7}
\end{equation*}

将求和次序颠倒后，可见它等于

\begin{equation*}
\frac{1}{2}\sum_{k,l} \left\{(\phi_{k}, \phi_{l}) + (\phi_{l}, \phi_{k})\right\}\left\{\Delta(s)_{ki}\Delta(s)_{lj} + \Delta(s)_{kj}\Delta(s)_{li}\right\}
\tag{2.6.8}
\end{equation*}

显然这个和属于 $[V \otimes V]$。定义在 $[V \otimes V]$ 上的 $\mathbf{S}$ 的表示记为 $[\Delta \boxtimes \Delta]$；同样，$\{V \otimes V\}$ 在 $\mathbf{S}$ 下不变，并给出一个表示 $\{\Delta \boxtimes \Delta\}$。这两个表示分别称为 $\Delta$ 的\textit{对称化克罗内克平方}与\textit{反对称化克罗内克平方}，有时分别记作 $[\Delta]^{2}$ 与 $\{\Delta\}^{2}$。

为计算对称化平方 $[\Delta \boxtimes \Delta]$ 的特征标，由式 (2.6.8) 得

\begin{equation*}
\chi_{[\Delta \boxtimes \Delta]}(s) = \frac{1}{2}\sum_{k,l} \left\{\Delta(s)_{kk}\Delta(s)_{ll} + \Delta(s)_{kl}\Delta(s)_{lk}\right\} = \frac{1}{2}\left\{\chi_{\Delta}^{2}(s) + \chi_{\Delta}(s^{2})\right\}.
\tag{2.6.9}
\end{equation*}

同样，对反对称化平方 $\{\Delta \boxtimes \Delta\}$ 即 $\{\Delta\}^{2}$：

\begin{equation*}
\chi_{\{\Delta \boxtimes \Delta\}}(s) = \frac{1}{2}\left\{\chi_{\Delta}^{2}(s) - \chi_{\Delta}(s^{2})\right\}.
\tag{2.6.10}
\end{equation*}

顺便指出，式 (2.6.10) 也可以由式 (2.6.6)、(2.6.9) 以及事实

\begin{equation*}
\Delta \boxtimes \Delta = [\Delta \boxtimes \Delta] + \{\Delta \boxtimes \Delta\}.
\tag{2.6.11}
\end{equation*}

推出。对有序对记号的含义感到略有困惑的读者，不妨把有序对的第一个成员看作属于粒子 1、第二个成员看作属于粒子 2，于是 $(\phi_{i}, \phi_{j}) - (\phi_{j}, \phi_{i})$ 这样的表达式就对应通常的反对称化波函数 $\phi_{i}(1)\phi_{j}(2) - \phi_{j}(1)\phi_{i}(2)$；式 (2.6.2) 的含义于此也是清楚的：算符 $s$ 同时、同样地作用于粒子 1 与粒子 2 两个空间。表示的对称化与反对称化平方的思想显然可以推广到群的表示的更高次乘积，例如见 Lyubarskii (1960) 第 4 章。

我们用晶体学点群提供的例子来说明式 (2.6.9) 与 (2.6.10) 在辨认 $\Delta$ 的对称化与反对称化乘积时的用法。对非简并的点群表示，式 (2.6.9) 与 (2.6.10) 可以简化：式 (2.6.9) 成为

\begin{equation*}
\chi_{[\Delta]^{2}}(s) = \{\chi_{\Delta}(s)\}^{2}
\tag{2.6.12}
\end{equation*}

而式 (2.6.10) 成为

\begin{equation*}
\chi_{\{\Delta\}^{2}}(s) = 0.
\tag{2.6.13}
\end{equation*}

换言之，对非简并的点群表示，对称化平方就是 $\Delta$ 的普通克罗内克平方，而反对称化平方为零。

对简并的点群表示，可以用式 (2.6.9)、(2.6.10) 与表 2.2 辨认 $[\Delta]^{2}$ 与 $\{\Delta\}^{2}$；例如对 $32$（$D_{3}$）的表示 $E$ 得

\begin{equation*}
[E]^{2} = A_{1} + E
\tag{2.6.14}
\end{equation*}

以及

\begin{equation*}
\{E\}^{2} = A_{2}.
\tag{2.6.15}
\end{equation*}

在表 2.8 中我们对各晶体学点群的简并表示给出了 $[\Delta]^{2}$ 与 $\{\Delta\}^{2}$；简并点群表示的对称化立方见 Cracknell and Joshua（1968）。

\begin{table}[H]
\centering
\textbf{表 2.8}\quad \textit{简并点群表示的} $[\Delta]^{2}$ \textit{与} $\{\Delta\}^{2}$

\smallskip
\begin{tabular}{|l|c|l|l|}
\hline
点群 & 表示 $\Delta$ & $[\Delta]^{2}$ & $\{\Delta\}^{2}$ \\ \hline
$\left.422\ (D_{4}),\ 4mm\ (C_{4v}),\ \bar{4}2m\ (D_{2d})\right\}$ & $E$ & $A_{1} + B_{1} + B_{2}$ & $A_{2}$ \\ \hline
$\left.32\ (D_{3}),\ 3m\ (C_{3v})\right\}$ & $E$ & $A_{1} + E$ & $A_{2}$ \\ \hline
$\left.622\ (D_{6}),\ 6mm\ (C_{6v})\right\}$ & $E_{1}$ & $A_{1} + E_{2}$ & $A_{2}$ \\ \hline
 & $E_{2}$ & $A_{1} + E_{2}$ & $A_{2}$ \\ \hline
$\bar{6}m2\ (C_{3h})$ & $E'$ & $A'_{1} + E'$ & $A'_{2}$ \\ \hline
 & $E''$ & $A''_{1} + E''$ & $A''_{2}$ \\ \hline
$23\ (T)$ & $T$ & $A + {}^{1}E + {}^{2}E + T$ & $T$ \\ \hline
$\left.432\ (O),\ \bar{4}3m\ (T_{d})\right\}$ & $E$ & $A_{1} + E$ & $A_{2}$ \\ \hline
 & $T_{1}$ & $A_{1} + E + T_{2}$ & $T_{1}$ \\ \hline
 & $T_{2}$ & $A_{1} + E + T_{2}$ & $T_{1}$ \\ \hline
\end{tabular}

\end{table}

\smallskip\noindent\textit{表 2.8 的注.} 本表转写自原书；其中 $23\ (T)$ 行 $[\Delta]^{2}$ 的具体组成请对照原书核对（原书印刷为 $A + {}^{1}E + {}^{2}E + T$ 的形式）。——中译者
